\section{问题三分析}

问题三将模型选择表述为场景效用与工作负载成本的双目标决策问题。对于模型 $i$ 和任务场景 $s$，Q2 负责提供场景效用 $U_i^{(s)}$；Q3 在此基础上引入场景化工作负载成本 $C_i^{(s)}$，并在预算约束下执行最优选型、Pareto 前沿分析、增量成本效益分析及稳健性检验。

\subsection{符号说明}

设 $i$ 表示模型，$s$ 表示任务场景，$N_s$ 表示场景标准工作负载中的调用次数，$L_{\mathrm{in}}^{(s)}$ 和 $L_{\mathrm{out}}^{(s)}$ 分别表示单次平均输入与计费输出 token 数，$p_{\mathrm{in},i}$ 与 $p_{\mathrm{out},i}$ 分别表示模型输入与输出 token 单价。这里的 $L_{\mathrm{out}}^{(s)}$ 指可计费输出 token（billable output tokens），而不只是用户可见回答 token；对于 reasoning/thinking 模型，若官方规则将内部推理或思考 token 计入输出计费，则该部分应纳入 $L_{\mathrm{out}}^{(s)}$。则场景成本为
\[
C_i^{(s)} = N_s\left(L_{\mathrm{in}}^{(s)}p_{\mathrm{in},i} + L_{\mathrm{out}}^{(s)}p_{\mathrm{out},i}\right).
\]
若价格以 USD/1M tokens 计，则在数值计算时需除以 $10^6$ 进行单位换算。

\subsection{工作负载成本模型}

Q3 不采用单一的 $Performance/Price$ 指标。原因在于 API 价格通常分别对输入与输出 token 计价，且不同任务场景的输入输出比例差异显著。因此成本必须视为任务工作负载的函数 $C_i^{(s)} = C_i(W_s)$，其中 $W_s=(N_s,L_{\mathrm{in}}^{(s)},L_{\mathrm{out}}^{(s)})$。

价格数据按模型逐项映射到官方 API 计费 SKU，并将模型身份映射、推理配置映射和价格状态分开记录。本阶段仅使用能够从官方来源对应到候选 API SKU 的价格；对于无法确认精确版本、推理配置或正式价格的模型，价格保留为 NA，不使用相邻型号或第三方报价替代。若官方价格以人民币计价，则按审计日 2026-08-17 的 USD/CNY 汇率转换为 USD/1M billable tokens，并同时保留原币种价格、汇率、汇率日期和汇率来源。所有自动采集或整理的价格在最终提交前仍需人工复核，因此不得标记为已人工验证。

三类基准工作负载分别为：科研长文本分析场景 $N_s=1, L_{\mathrm{in}}=16000, L_{\mathrm{out}}=4000$；通用日常任务场景 $N_s=1, L_{\mathrm{in}}=2000, L_{\mathrm{out}}=1000$；代码开发场景 $N_s=1, L_{\mathrm{in}}=6000, L_{\mathrm{out}}=6000$。这些工作负载是用于模型间可比性的标准化情景设定，并非行业平均调用量。由于 $N_s$ 只会在同一场景内线性放大所有模型成本，设定 $N_s=1$ 可直接比较单次标准任务下的成本排序。

\subsection{Pareto 双目标决策模型}

Q3 的核心优化目标为
\[
\max U_i^{(s)}, \qquad \min C_i^{(s)}.
\]
模型 $j$ 支配模型 $i$ 当且仅当
\[
U_j^{(s)} \ge U_i^{(s)}, \qquad C_j^{(s)} \le C_i^{(s)},
\]
且至少一项严格不等。未观测到的效用或成本不得自动视为 0。

\subsection{预算约束选型模型}

给定预算 $B$，最优模型定义为
\[
i^*(B,s)=\arg\max_i U_i^{(s)} \quad \text{s.t.}\quad C_i^{(s)}\le B.
\]
当预算不足以支持任何模型时，输出无可行解标记。预算从低到高变化时，可得到模型切换点序列，用于描述不同预算区间下的最优选型规则。

\subsection{增量成本效益模型}

在 Pareto 有效前沿上，按成本从低到高排序后，定义相邻模型之间的增量成本效益比为
\[
ICER_k = \frac{C_{k+1}-C_k}{U_{k+1}-U_k}.
\]
当 $\Delta U=0$ 时需单独处理，不得简单除零。该指标用于解释额外一单位场景效用对应的边际成本。

\subsection{成本--性能规律拟合}

Q3 对全体模型与 Pareto 前沿分别拟合三类关系：线性模型 $U=a+bC$、对数模型 $U=a+b\ln(1+C)$、饱和模型 $U=a-be^{-kC}$。每个场景分别报告 $R^2$、AIC、AICc、BIC 与 LOOCV error。鉴于核心模型数约为 10，AICc 必须保留。该分析属于横截面相关分析，不构成因果推断。

\subsection{敏感性与稳健性分析}

Q3 设置三类敏感性分析：输入 token 规模扰动、输入/输出比例扰动、价格扰动。分析对象包括成本排序、Pareto 成员、预算切换点以及最优模型变化。若未来获得 Q2 的 bootstrap 场景效用样本，还可进一步估计模型进入 Pareto 有效前沿的概率。

输入规模敏感性采用基准输入长度的 $0.5\times,1.0\times,2.0\times$，并保持输出长度不变；输入/输出比例敏感性则在固定输入长度下将输出长度设为输入长度的 $0.1,0.3,0.5,1.0,2.0$ 倍。价格敏感性对输入和输出单价同时施加 $-10\%,-5\%,0,+5\%,+10\%$ 扰动。对于存在峰谷时段价格、缓存折扣或批处理折扣的官方价格，主分析采用标准实时 API 的未缓存价格，折扣价格仅作为敏感性或补充分析变量。若 Q2 的正式场景效用文件尚未提供，Q3 仅输出成本与数据准备状态，不生成最终 Pareto 排名或模型推荐。

对于带 fallback 的 benchmark 配置，Q3 不将基础模型价格自动等同于完整配置成本。若冻结记录显示实际未触发 fallback，则可使用基础模型价格；若发生 fallback 且目标模型与 token 记录可追溯，则按真实记录进行成本修正；若仅知道发生或允许 fallback 但缺少目标模型或 token 轨迹，则该模型价格状态标记为特殊案例，不进入最终全模型成本比较。

\subsection{缺失能耗数据处理原则}

对于闭源 API 模型，若官方未公开单次调用能耗或可靠推理硬件信息，则应记为 NA，不得使用参数量、GPU 功耗或吞吐量自行推算能耗，也不得将缺失值视为 0。
